Recent studies on AI coding agents have proven the efficacy of using Large Language Models (LLMs) to handle coding tasks such as code generation and debugging. To apply LLMs in daily software engineering scenarios, both industry and academia have proposed many agentic frameworks. Industry-proprietary examples of agentic frameworks include Google's Colab AI \cite{google_2025_colab_ai_first}, Claude Code \cite{anthropic_2025_claude_code_ga}, and Cursor Agent \cite{cursor_2025_agent_cli}. There is a continuous and rapid growth of user bases for the above frameworks, reflecting their perceived utility in real-world software engineering tasks. This implies increasing demand for production-ready AI assistants that integrate seamlessly into existing development workflows. Besides those proprietorial frameworks, there are also many open source implementations from both academia and industry. Examples of these open-source agentic frameworks include Gemini Cli \cite{mullen_salva_2025_gemini_cli}, SWE-Agent \cite{yang2024sweagent}, AutoCodeRover \cite{zhang2024autocoderover}, Agentless \cite{xia2024agentlessdemystifyingllmbasedsoftware}, etc. These frameworks enable existing LLMs to achieve human-level performance when prompted with clear instructions and sufficiently rich context. Despite these advances, current agentic systems still face two major challenges: 1. How to provide accurate and complete context information to the LLM for the best performance / code quality (later referred to as \textbf{Challenge 1}) \cite{ouyang2025repographenhancingaisoftware}; 2. How to ensure fast patch generation when the LLM is fed with a large number of requests and / or a large number of tokens during codebase exploration (later referred to as \textbf{Challenge 2}). 

More recently, there have been studies that address Challenge 1 or Challenge 2. However, no existing system attempts to tackle both challenges at once. There are many reasons why joint optimization is highly desirable, the most apparent being cost. For a complex code repository and an intricate issue, it often takes a significantly long time for the agent to generate a patch, even when the patch is unlikely to resolve the issue given the complexity. And if the patch is deemed problematic, the consumed time and token length all contribute to the cumulative engineering expense.

This thesis project aims to fill the gap in understanding: what are the solutions used by existing frameworks to address the challenges, how these solutions affect the performance, and what trade-offs need to be considered for joint optimization of context retrieval and patch generation latency. Our results and analysis illustrate the possibility as well as difficulty of integrating such solutions to step towards a fast yet performant code LLM-based agentic framework. The paper is organized as follows:

\mytodo{update below once restructured}
\vspace{0.1in}
\begin{itemize}
    \item In Chapter~\ref{cha:ch1}, we inspect the presence of Challenge 1 in existing systems, then evaluate the performance of RepoGraph \cite{ouyang2025repographenhancingaisoftware}, a state-of-the-art solution for Challenge 1. The evaluation primarily concerns consistency of performance using different model backends.
    \item In Chapter~\ref{cha:ch2}, we inspect the presence of Challenge 2 in existing systems, and evaluate the latency of CacheBlend \cite{yao2025cacheblendfastlargelanguage}, a state-of-the-art solution targeting Challenge 2. The evaluation involves a comparison with legacy systems such as vLLM and SGLang \cite{kwon2023vllm, zheng2024sglangefficientexecutionstructured}. The primary distinction of our setup from that of the authors is that this evaluation uses RepoGraph produced agent prompts as dataset, and is meant to measure latency under real-world code agent use case scenarios. 
    \item In Chapter~\ref{cha:ch3}, we discuss how the above solutions can be integrated into a unified end-to-end code agent system that is reliable, fast, and provides high-quality code generation results. 
\end{itemize}